\subsection{Test di intergrazione}
\label{subsec:test_integrazione}
Lo scopo dei test di integrazione è individuare eventuali difetti di progettazione o lacune nei test unitari, 
verificando che i diversi moduli o componenti del sistema interagiscano correttamente tra loro quando vengono integrati.
\begin{tabularx}{\textwidth}{|c|X|c|}
\hline
\textbf{ID} & \textbf{Descrizione} & \textbf{Stato}\\
\hline
T.I.1 & Verificare che il componente "StandardPage" soddisfi le seguenti condizioni:
\begin{itemize}%
    \item deve restituire correttamente la lista dei Dataset salvati dall'interno del database tramite DatasetService;% sicuramente test di integrazione di back-end
    \item deve restituire correttamente la lista dei Test salvati dall'interno del database tramite TestService;
    \item deve restituire correttamente la lista dei LLM salvati dall'interno del database tramite LLMService;
\end{itemize}
& -\\
\hline
T.I.2 & Verificare che il componente "HomePage" soddisfi le seguenti condizioni:
\begin{itemize}
    \item 
\end{itemize}
& -\\
\hline
T.I.3 & Verificare che il componente "DatasetListPage" soddisfi le seguenti condizioni:
\begin{itemize}
    \item Deve restituire correttamente i risultati della ricerca sotto forma di elementi della lista  di dataset filtrati quando si effettua una ricerca;
    \item l'utente viene reindirizzato nella schermata apposita dopo aver premuto il pulsante di caricamento di un dataset a partire da un file JSON;
    \item l'utente viene reindirizzato alla pagina del dataset caricato dopo aver premuto il pulsante di caricamento a schermo del dataset;
    \item l'utente viene reindirizzato nella schermata apposita dopo aver premuto il pulsante di rinominazione del dataset;
    \item nel caso in cui l'operazione di caricamento dei dati vada a buon fine, viene visualizzato un messaggio di successo;
    \item nel caso in cui l'operazione di caricamento dei dati fallisse, viene visualizzato un messaggio di errore;
\end{itemize}
& -\\
\hline
T.I.4 & Verificare che il componente "DatasetContentPage" soddisfi le seguenti condizioni:
\begin{itemize}
    \item deve restituire correttamente la lista delle coppie di domande e risposte dall'interno del database tramite QAService;
    \item L'utente viene reindirizzato alla pagina della selezione degli LLM quando si preme il pulsante "Esegui test";
    \item viene visualizzato correttamente il form che permette l'inserimento della coppia di domanda e risposta;
    \item nel caso in cui l'operazione di ottenimento dei dati fallisse, viene visualizzato un messaggio di errore;
    \item nel caso in cui l'operazione di salvataggio dei dati vada a buon fine, viene visualizzato un messaggio di successo;
    \item nel caso in cui l'operazione di esecuzione del test fallisse a causa di elementi incompleti, viene visualizzato un messaggio di errore;
\end{itemize}
& -\\
\hline
T.I.5 & Verificare che il componente "TestPage" soddisfi le seguenti condizioni:
\begin{itemize}
    \item nel caso in cui l'operazione di ottenimento dei dati fallisse, viene visualizzato un messaggio di errore;
    \item nel caso in cui l'operazione di salvataggio del test temporaneo fallisse, viene visualizzato un messaggio di errore;
\end{itemize}
& -\\
\hline
T.I.6 & Verificare che il componente "TestListPage" soddisfi le seguenti condizioni:
\begin{itemize}
    \item l'utente viene reindirizzato nella schermata apposita dopo aver premuto il pulsante di rinominazione del test;
\end{itemize}
& -\\
\hline
T.I.7 & Verificare che il componente "TestComparisonPage" soddisfi le seguenti condizioni:
\begin{itemize}
    \item nel caso in cui l'operazione di ottenimento dei dati fallisse, viene visualizzato un messaggio di errore;
\end{itemize}
& -\\
\hline
T.I.8 & Verificare che il componente "LLMListPage" soddisfi le seguenti condizioni:
\begin{itemize}
    \item 
\end{itemize}
& -\\
\hline
T.I.9 & Verificare che il componente "LLMPage" soddisfi le seguenti condizioni:
\begin{itemize}
    \item 
\end{itemize}
& -\\
\hline
\end{tabularx}
